\documentclass[a4paper]{report}
% Shared preamble for the English probability and statistics documents.
% Based on the original course preamble. Compile with pdfLaTeX.
\usepackage[top=1.5cm,bottom=2cm,left=2cm,right=2cm]{geometry}
\usepackage[T1]{fontenc}
\usepackage{lmodern}
\usepackage[utf8]{inputenc}
\usepackage[english]{babel}

\usepackage{amsmath}
\usepackage{amssymb}
\usepackage{amsthm}
\usepackage{mathrsfs}
\usepackage{mathtools}
\usepackage{etoolbox}
\usepackage{tikz}
\usetikzlibrary{shapes.multipart,arrows.meta}
\usepackage{xcolor}
\usepackage{enumerate}
\usepackage[shortlabels]{enumitem}
\usepackage{hyperref}
\usepackage{version}
\usepackage{bbold}
\usepackage{fontawesome5}
\usepackage{fancyhdr}
\lfoot{}
\usepackage{natbib}
\usepackage{graphicx}
\usepackage{pdfpages}
\usepackage{authblk}
    
\usepackage{blindtext}
\usepackage{multicol}
\usepackage{collcell}
\usepackage{wrapfig}
\usepackage{blkarray}
\usepackage{cancel}
\usepackage{appendix}
\usepackage{float}
\makeatletter
\input{tex/assets/slashbox.sty}
\makeatother
\hypersetup{unicode,colorlinks,citecolor=red,filecolor=black,
  linkcolor=red,urlcolor=cyan}
\setlength{\columnsep}{1cm}
\pagenumbering{gobble}

\def\numtext#1{%
  \ifcase#1\or one\or two\or three\or four\or five\or six\or seven\or eight\or nine\or ten\or
  eleven\or twelve\or thirteen\or fourteen\or fifteen\or sixteen\or seventeen\or eighteen\or nineteen\or twenty\or Lots\fi}

\newcommand{\Byte}[1]{%
  \def\nodecontents{}%
  \foreach[count=\l] \k in {#1} {\xappto{\nodecontents}{\noexpand\nodepart{\numtext{\l}}\k}}%
  \node[rectangle split,rectangle split parts=8,rectangle split horizontal,draw] at (2,2){\nodecontents};}

\newcommand{\D}{{\mathrm{d}}}
\newcommand{\cE}{{\mathcal{E}}}
\newcommand{\cG}{{\mathcal{G}}}
\newcommand{\cX}{{\mathcal{X}}}
\newcommand{\cY}{{\mathcal{Y}}}
\newcommand{\cB}{{\mathcal{B}}}
\newcommand{\cF}{{\mathcal{F}}}
\newcommand{\cO}{{\mathcal{O}}}
\newcommand{\cS}{{\mathcal{S}}}
\newcommand{\cT}{{\mathcal{T}}}
\newcommand{\given}{\, \big| \,}
\newcommand{\cL}{{\mathscr{L}}}
\newcommand{\cD}{{\mathcal{D}}}
\newcommand{\cA}{{\mathscr{A}}}
\newcommand{\cP}{{\mathscr{P}}}
\newcommand{\cR}{{\mathscr{R}}}
\newcommand{\A}{{\mathcal A}}
\newcommand{\one}{\mathbb{1}}
\newcommand{\mathbbm}[1]{\mathbb{#1}}

\providecommand{\C}{}
\renewcommand{\C}{\mathbb C}
\newcommand{\R}{\mathbb R}
\newcommand{\Q}{\mathbb Q}
\newcommand{\N}{\mathbb N}
\newcommand{\T}{\mathbb T}
\newcommand{\F}{\mathbb F}
\renewcommand{\P}{\mathbb P}
\newcommand{\Z}{\mathbb Z}
\newcommand{\E}{\mathbb {E}}
\newcommand{\Pbb}[1]{\P\left( #1\right)}
\newcommand{\Ebb}[1]{\E\left[ #1\right]}

\newcommand{\lb}{\left (}
\newcommand{\rb}{\right )}
\newcommand{\lbb}{\left [}
\newcommand{\rbb}{\right ]}
\newcommand{\labs}{\left |}
\newcommand{\rabs}{\right |}
\newcommand{\lbrb}[1]{\lb #1 \rb}
\newcommand{\lbbrbb}[1]{\lbb#1\rbb}
\newcommand{\labsrabs}[1]{\labs#1\rabs}
\newcommand{\lbbrb}[1]{\lbb#1\rb}
\newcommand{\lbrbb}[1]{\lb#1\rbb}
\newcommand{\langlerangle}[1]{\langle#1\rangle}
\newcommand{\lbcurly}{\{}
\newcommand{\rbcurly}{\}}
\newcommand{\curly}[1]{\lbcurly#1\rbcurly}
\newcommand{\abs}{\labsrabs}
\newcommand{\ind}[1]{1_{#1}}
\newcommand{\indep}{\perp \!\!\! \perp}
\renewcommand{\d}{\, \mathrm{d}}

\newtheorem{theorem}{Theorem}
\newtheorem{problemCnt}{Lemma}
\theoremstyle{definition}
\newtheorem{problem}[problemCnt]{Problem}
\newtheorem{problembook}[problemCnt]{\faBook \,  Problem}
\newtheorem{problempencil}[problemCnt]{\faPencil* \,  Problem}
\newtheorem{problemstar}[problemCnt]{\faStar\,  Problem}
\newtheorem{problemstarpencil}[problemCnt]{\faStar \, \faPencil* \,   Problem}
\newtheorem{probleminteresting}[problemCnt]{\faSmile[regular] \,   Problem}
\newtheorem{solution}{Solution}

\newcommand{\fwcell}[1]{\makebox[\arraystretch\normalbaselineskip]{$#1$}}
\newenvironment{enumeratei}{\begin{enumerate}[(i)]}{\end{enumerate}}
\newenvironment{enumeratea}{\begin{enumerate}[(a)]}{\end{enumerate}}
\newenvironment{enumeraten}{\begin{enumerate}[\upshape 1.]}{\end{enumerate}}

% Compatibility helpers needed by a few files in the combined archive.
\AtBeginDocument{\noindent{\footnotesize English translation prepared with ChatGPT (OpenAI, GPT-6), 26 October 2026, for translation of the Bulgarian text into English and LaTeX formatting.
The Bulgarian originals remain the reference; translation may contain errors.}\par\smallskip}
\newcommand{\MissingGraphic}[1]{%
  \begin{center}\fbox{\parbox{0.82\linewidth}{\centering Missing source image: \texttt{#1}}}\end{center}}



\setlength{\emergencystretch}{3em}
\allowdisplaybreaks
\begin{document}


\noindent Probability and Statistics, Software Engineering, Year 3, Groups 4/5, 2020/2021 \hfill Test 2

\hrulefill

\begin{it}
Your grade is 2 plus the number of points earned. Time allowed: 4 hours. Good luck.

Throughout, we work on a probability space $(\Omega, \cF, \P)$.

\end{it}

\begin{problem} (2 points)
\begin{enumerate}
    \item (0.5 points) Two cards are drawn from a standard 52-card deck. Let $X$ be the number of aces and $Y$ the number of hearts. Find the joint distribution and correlation of $X,Y$.

$\bullet$ Let $C=\binom{52}{2}=1326$. Combinatorial counting gives
\begin{center}
\begin{tabular}{ |c|c|c|c| }
 \hline
$X\ (\text{aces})\backslash Y\ (\text{hearts})$ & 0 & 1 & 2 \\
  \hline
 0 & $\binom{36}{2}/ C$& $\binom{12}{1} \binom{36}{1}/ C$ &   $\binom{12}{2}/ C$\\
 1 &  $\binom{3}{1} \binom{36}{1}/ C$ &  ($\binom{12}{1} \binom{3}{1} +
 \binom{1}{1} \binom{36}{1})/ C$ &  $\binom{12}{1}/ C$\\
 2 & $\binom{3}{2}/ C$& $\binom{3}{1}/ C $& 0\\
 \hline\end{tabular} \, ,
 \end{center}
After calculating and summing the rows and columns (or directly by counting),
 \begin{center}
\begin{tabular}{ |c|c|c|c| }
 \hline
$X\ (\text{aces})\backslash Y\ (\text{hearts})$ & 0 & 1 & 2 \\
  \hline
 0 & $105/221$& $72/221$ &   $11/221$\\
 1 &  $18/221$ &  $12/221$ &  $2/221$\\
 2 & $1/442$& $1/442$ & 0\\
 \hline
\end{tabular}
\\
\begin{tabular}{ |c|c|c|c| }
 \hline
$Y\ (\text{hearts})$ & 0 & 1 & 2 \\
  \hline
 $\P(Y = \cdot) $& $19/34$& $13/34$ &   $2/34$\\
 \hline
\end{tabular}
\\
\begin{tabular}{ |c|c|c|c| }
 \hline
$X\ (\text{aces})$ & 0 & 1 & 2 \\
  \hline
 $\P(X = \cdot) $& $188/221$& $32/221$ &   $1/221$\\
 \hline
\end{tabular} \, .

\end{center}
Therefore
\begin{align*}
    &\E X = 34/221 = 2/13;    &D X = \E X^2 - (\E X)^2 = 36/221 - (2/13)^2 =
    400/2873;\\
    &\E Y = 17/34 = 1/2;    &D Y =  \E Y^2 - (\E Y)^2 = 21/34 - 1/4 = 25/68;\\
    &\E XY = 12/221 + 4/221 + 2/442 = 1/13;
    &Cor(X,Y) = \frac{\E XY - \E X \E Y}{\sqrt{D X \cdot D Y}} =
    \frac{1/13 - 1/13 }{\sqrt{D X \cdot D Y}} = 0
    \, .
\end{align*}
Thus $X,Y$ are uncorrelated, although they are not independent!



\item (1 point)
Let $X$ be an asset price and $Y$ the average deposit interest rate (both daily and in dimensionless units). Their joint density is
$f_{X,Y}(x,y)=c(x^2+e^yx)$ for $0<x,y<1$, and zero elsewhere, where $c$ is
a constant. Find
$c$ and the covariance of $X,Y$. What is the expected asset price if
the interest rate
is $0.5$?

$\bullet$
\[
1 = c \int_0^1 \int_0^1 (x^2 + e^y x) \d x \d y = c \int_0^1 \frac13 + \frac12 e^y \d y
= c\left(\frac13 + \frac12 ( e - 1)\right) = \frac{3e - 1}{6} \, .
\]
Hence $c=6/(3e-1)$. For $x\in(0,1)$,

\begin{align*}
f_X(x) = \int_0^1 c(x^2 + e^yx) \d y = c(x^2 + (e - 1)x) \, .
\end{align*}
Therefore
\[ \E X = c\int_0^1 x ( x^2 + (e-1)x) \d x = c\left(\frac14 + (e-1)\frac13\right)
= c\frac{4e - 1}{12}\, .\]
Similarly, for $y\in(0,1)$,
\[
f_Y(y) = c\int_0^1 x^2 + e^y x \d x = c\left(\frac13 + \frac12 e^y\right)
= \frac{c}{6}\left(2 + 3 e^y\right)\, ,
\]
\[
\E Y = \frac{c}{6}\int_0^1 y (2 + 3 e^y) \d y  =
\frac{c}{6}\left(
2 \frac12 + 3 \right) = \frac{2c}{3} \, ,
\]
where we used
$\int xe^x\d x=\int x\d e^x=xe^x-\int e^x\d x=e^x(x-1)$ and
hence $\int_0^1 xe^x\d x=1$.
Next,
\[
\E XY = c \int_0^1 \int_0^1 xy(x^2 + e^y x) \d x \d y
=c \int_0^1 \frac14 y + \frac13 ye^y \d y =
c \left( \frac14 \cdot\frac12 + \frac13 \right) = \frac{11c}{24}\, .
\]
Therefore
\[
Cov(X, Y) = \E XY - \E X \E Y = \frac{11c}{24} - c^2 \frac{4e - 1}{18} \approx
-0.0014 \, .
\]
For the last question,
\begin{align*}
\E (X | Y = 0.5) &= \int_0^1 x f_{X|Y}(x | y = 0.5) \d x
= \int_0^1 x \frac{f_{X, Y}(x, 0.5)}{f_Y(0.5)} \d x
= \int_0 ^1 x \frac{c(x^2 + x \sqrt{e})}{\frac{c}{6}\left(2 + 3 \sqrt{e}\right)}\d x \\
&= \frac{6}{2 + 3\sqrt{e}} \int_0^1 x^3 + x^2 \sqrt{e} \d x
= \frac{6}{2 + 3\sqrt{e}} (\frac14 + \frac13 \sqrt{e}) \approx 0.6907 \, .
\end{align*}


\item (0.5 points) Let $X$ have density $f_X(x)=\frac14xe^{-x/2}$ for $x\geq0$ and
zero otherwise. Find the density of $Y=-2X+2$.

$\bullet$ Apply the transformation theorem directly or follow its proof:
for $y>2$, $\P(Y\leq2)=\P(X\geq0)=1$ and $f_Y(y)=0$.

For $y\leq2$,
$F_Y(y)  = \P(Y \leq y) = \P(-2X + 2 \leq y) = \P( X \geq 1/2(2 - y)) = 1 - F_X( 1/2(2 - y))$.
Therefore \begin{align*}
f_Y(y) = \frac {\d}{\d y}F_y(y) &= \frac {\d}{\d y}\left(1 - F_X\left( \frac12\left(2 - y\right)
\right)\right)
= -F_X'\left( \frac12\left(2 - y\right)\right) \cdot \left( \frac12\left(2 - y\right)\right) '
\\
&= \frac12 f_X\left(1 - \frac12 y\right) = \frac12 \cdot \frac14 \left(1- \frac12y\right)e^{y/4 - 1/2} \, .
\end{align*}
    \end{enumerate}
\end{problem}


\begin{problem}(0.5 points)

According to official data, Sofia's historical average January temperature is $-0.5^\circ C$. Let $X$ be the average January temperature in Sofia in degrees Celsius, assumed normally distributed. The variance is usually estimated
from data too, but here we consider a different (not recommended) method.

Suppose someone regards a temperature above $20^\circ C$ as an extreme anomaly
with probability 1 percent. What variance estimate does this imply? What
is then the probability of a temperature below $-15^\circ C$?

$\bullet$
Let $\sigma^2$, $\sigma\geq0$, be the variance, so $X\sim N(-0.5,\sigma^2)$. We have
\[\frac{1}{100} = \P\left(X \geq 20\right) = \P\left( \frac{X + 0.5}{\sigma} \geq \frac{20.5}{\sigma}\right)
= \P\left(N(0,1) \geq \frac{20.5}{\sigma}\right) = 1 - \Phi\left(\frac{20.5}{\sigma}\right) \, ,
\]
where $\Phi(x)=\P(N(0,1)\leq x)$. From a normal table, find $z$
such
that $\Phi(z)=0.99$. This gives $20.5/\sigma\approx2.33$ and
$\sigma\approx8.80$. Then
\[
\P(X \leq -15) =\P\left( \frac{X + 0.5}{8.80} \leq \frac{-14.5}{8.80}\right)
\approx \P\left(N(0,1) \leq -1.65\right) \approx 0.0495 \, .
\]

\end{problem}

\begin{problem}(1.5 points)
A point $A$ is chosen randomly in the disk $k(O,1)$ with centre $O$ and radius $1$. Let
$X=|OA|$. Can you guess
the mode and median? Justify your guesses. What would you expect $\E X$ to be?
\textit{(The mode of a discrete distribution is a value with maximum probability.
In the continuous case, analogously, we seek a value maximizing $f_X$.
We call $a$ a median of $X$ if $\P(X\leq a)=\P(X\geq a)=1/2$.)}

$\bullet$ Informally, we may guess that the mode is 1, because at a fixed
distance from the centre, a larger distance corresponds to more points.
Formally, the cardinalities are of course the same. The median is the value
satisfying $\pi r^2=1-\pi r^2$.
    \begin{enumerate}

    \item (0.75 points) Find the distribution function, density, expectation and
    variance of $X$.

$\bullet$ For $c\in(0,1)$, $F_X(c)=\P(X\leq c)=\P(X\in k(O,c))=S_{k(O,c)}/S_{k(0,1)}
=c^2$. Hence $f_X(c)=F_X'(c)=2c$. For completeness, if $c\geq1$, then $F_X(c)=1,
f_X(c)=0$, while if $c\leq0$, $F_X(c)=0$, $f_X(c)=0$.
We obtain
\[\E X = \int_0^1 x 2x \d x = \frac23; D X = \int_0^1 x^2 \cdot 2x \d x - \frac49= \frac{1}{18}\, . \]
    \item (0.75 points) Now let $A_1,A_2,A_3$ be three points chosen randomly
    and independently in the same disk. What is the expected distance to the point
    nearest the centre? And to the farthest? \textit{(Bonus: Find the expected
    distance to the middle point. Should it equal $\E X$?)}

$\bullet$
Let the distances from $O$ to $A_1,A_2,A_3$ be $X_1,X_2,X_3$, respectively. Then
$X_1,X_2,X_3\sim X$ are independent, where $X$ is as in the previous part.
The distance to the nearest point is $Y_1=\min(X_1,X_2,X_3)$, and to
the farthest, $Y_2=\max(X_1,X_2,X_3)$.
For $y\in(0,1)$,
$\P(Y_2\leq y)=\P(X\leq y)^3=y^6$. Therefore
\[
f_{Y_2}(y) = 6y^5; \E Y_2 = \int_0^1 y 6y^5 \d y = \frac67 \, .
\]
Similarly, $\P(Y_1\leq y)=1-\P(Y_1\geq y)=1-(\P(X\geq y))^3=
1-(1-y^2)^3$. Therefore
\[
f_{Y_1}(y) = 6y(1-y^2)^2; \E Y = \int_0^1 y 6y(1 - y^2)^2 \d y = \frac{16}{35} \, .
\]
\end{enumerate}
\end{problem}

\end{document}
